% ============================================================================
%  CLIPP — MILP Formulation (academic style)
% ============================================================================
\documentclass[11pt]{article}
\usepackage[margin=2.5cm]{geometry}
\usepackage{amsmath,amssymb}
\usepackage{enumitem}
\usepackage{booktabs}

\newcommand{\setN}{\mathcal{N}}
\newcommand{\setJ}{\mathcal{J}}
\newcommand{\Jm}{\mathcal{J}_M}
\newcommand{\Jo}{\mathcal{J}_O}
\newcommand{\Jc}{\mathcal{J}_C}
\newcommand{\setA}{\mathcal{A}}
\newcommand{\Aj}{A_j}
\newcommand{\setK}{\mathcal{K}}
\newcommand{\depot}{S}

\title{Climate-Integrated Street Cleaning Routing Problem\\(MILP Formulation)}
\author{}
\date{}

\begin{document}
\maketitle

\section{Notation}

Let $\mathcal{N} = \{0, \dots, N-1\}$ be the set of junctions and
$\mathcal{J} = \{0, \dots, M-1\}$ the set of streets in the city network.
The streets are partitioned into three categories:
$\mathcal{J}_M$ (mandatory),
$\mathcal{J}_O$ (optional), and
$\mathcal{J}_C$ (connectors), so that
$\mathcal{J} = \mathcal{J}_M \cup \mathcal{J}_O \cup \mathcal{J}_C$.
Let $\mathcal{A}$ be the set of directed arcs obtained from the street
network, where each two-way street generates two arcs and each one-way
street generates one.  For a street $j \in \mathcal{J}$, denote by
$A_j \subseteq \mathcal{A}$ the set of arcs corresponding to~$j$.  For a
junction $u \in \mathcal{N}$, let $\delta^+(u)$ and $\delta^-(u)$ be the
sets of arcs leaving and entering~$u$, respectively.

Let $\mathcal{K} = \{0, \dots, C-1\}$ be the fleet of cleaning vehicles.
For each arc $a \in \mathcal{A}$, denote by $c_a$ its traversal time in
seconds.  For each street $j \in \mathcal{J}$, denote by $\ell_j$ its
length in metres and by $r_j$ its cleaning requirement in litres per
kilometre.  For each vehicle $k \in \mathcal{K}$, denote by $q_k$ its water
capacity in litres per kilometre.  A vehicle~$k$ can clean street~$j$ only
if $q_k \ge r_j$; connector streets ($j \in \mathcal{J}_C$) cannot be
cleaned.

Let $T$ be the shared time limit for every vehicle, $S \in \mathcal{N}$ the
depot junction at which all vehicles start and must return, and $W \ge 0$
the water-waste penalty factor.  Define $U_a = \min(\lfloor T / c_a
\rfloor,\, 50)$ as the traversal upper bound for arc~$a$ and $\mathcal{M}
= 1 + \sum_{a \in \mathcal{A}} U_a$ as the Big-M constant.

Denote by $f_{a,k}$, $a \in \mathcal{A}$, $k \in \mathcal{K}$, the
integer decision variable representing the number of times vehicle~$k$
traverses arc~$a$.  Denote by $x_{j,k}$, $j \in \mathcal{J} \setminus
\mathcal{J}_C$, $k \in \mathcal{K}$, the binary decision variable which
takes the value~1 if vehicle~$k$ cleans street~$j$, and~0 otherwise.  If
$q_k < r_j$, the variable $x_{j,k}$ is fixed to~0.  Finally, denote by
$y_j$, $j \in \mathcal{J} \setminus \mathcal{J}_C$, the binary decision
variable which takes the value~1 if street~$j$ is cleaned by any vehicle,
and~0 otherwise.


\begin{center}
\begin{tabular}{@{} l l @{}}
\toprule
Symbol & Description \\
\midrule
$\setN = \{0, \dots, N-1\}$       & Junctions (nodes) \\
$\setJ = \{0, \dots, M-1\}$       & Streets (edges) \\
$\Jm \subset \setJ$               & Mandatory streets \\
$\Jo \subset \setJ$               & Optional streets \\
$\Jc \subset \setJ$               & Connector streets ($\setJ = \Jm \cup \Jo \cup \Jc$) \\
$\setA$                            & Directed arcs (a two-way street generates two arcs) \\
$\Aj \subseteq \setA$             & Set of arcs corresponding to street $j$ \\
$\setK = \{0, \dots, C-1\}$       & Vehicles \\
$\delta^+(u), \delta^-(u)$         & Outgoing / incoming arcs at node $u$ \\
\bottomrule
\end{tabular}
\end{center}


\begin{center}
\begin{tabular}{@{} l l @{}}
\toprule
Symbol & Description \\
\midrule
$T$            & Shared time limit for every vehicle (seconds) \\
$\depot$       & Depot junction \\
$W$            & Water-waste penalty factor \\
$c_a$          & Traversal time of arc $a$ (seconds) \\
$\ell_j$       & Length of street $j$ (metres) \\
$r_j$          & Cleaning requirement of street $j$ (litres\,/\,km) \\
$q_k$          & Water capacity of vehicle $k$ (litres\,/\,km) \\
$U_a$          & Traversal upper bound for arc $a$: $U_a = \min\!\bigl(\lfloor T / c_a \rfloor,\; 50\bigr)$ \\
$\mathcal{M}$  & Big-M constant: $\mathcal{M} = 1 + \sum_{a \in \setA} U_a$ \\
\bottomrule
\end{tabular}
\end{center}


\section{Decision variables}
%% -------------------------------------------------------------------

\begin{align}
f_{a,k} &\in \mathbb{Z}_{\ge 0}
    &&\forall\, a \in \setA,\; k \in \setK
    \tag{flow variable}\\[4pt]
x_{j,k} &\in \{0,1\}
    &&\forall\, j \in \setJ \setminus \Jc,\; k \in \setK
    \tag{cleaning assignment}\\[4pt]
y_j     &\in \{0,1\}
    &&\forall\, j \in \setJ \setminus \Jc
    \tag{cleaned indicator}
\end{align}

\begin{itemize}[nosep]
  \item $f_{a,k}$ counts how many times vehicle~$k$ traverses arc~$a$.
  \item $x_{j,k} = 1$ if vehicle~$k$ cleans street~$j$.
  \item $y_j = 1$ if street~$j$ is cleaned by \emph{any} vehicle.
\end{itemize}

\medskip\noindent
\textbf{Note.}
If vehicle~$k$ cannot clean street~$j$ (i.e.\ $q_k < r_j$, or $j \in \Jc$),
$x_{j,k}$ is fixed to~$0$ and not created as a decision variable.

\section{Formulation}

Using the above definitions, the Climate-Integrated Street Cleaning Routing
Problem can be formulated as follows:

\begin{align}
\text{Maximise} \quad
& \sum_{j \in \mathcal{J} \setminus \mathcal{J}_C} \ell_j \, y_j
\;-\;
W \sum_{j \in \mathcal{J} \setminus \mathcal{J}_C}
  \sum_{k \in \mathcal{K}}
  \frac{(q_k - r_j)\,\ell_j}{1000}\, x_{j,k}
&& \tag{1} \label{eq:obj} \\[10pt]
%
\text{subject to} \notag \\[4pt]
%
& \sum_{a \in \delta^+(u)} f_{a,k}
  \;=\;
  \sum_{a \in \delta^-(u)} f_{a,k}
&& \forall\, u \in \mathcal{N},\; k \in \mathcal{K},
\tag{2} \label{eq:flow} \\[6pt]
%
& \sum_{a \in \mathcal{A}} c_a \, f_{a,k} \;\le\; T
&& \forall\, k \in \mathcal{K},
\tag{3} \label{eq:time} \\[6pt]
%
& \sum_{a \in A_j} f_{a,k} \;\ge\; x_{j,k}
&& \forall\, j \in \mathcal{J} \setminus \mathcal{J}_C,\;
   k \in \mathcal{K},
\tag{4} \label{eq:traverse} \\[6pt]
%
& x_{j,k} \;\le\; y_j
&& \forall\, j \in \mathcal{J} \setminus \mathcal{J}_C,\;
   k \in \mathcal{K},
\tag{5} \label{eq:link-fwd} \\[6pt]
%
& y_j \;\le\; \sum_{k \in \mathcal{K}} x_{j,k}
&& \forall\, j \in \mathcal{J} \setminus \mathcal{J}_C,
\tag{6} \label{eq:link-rev} \\[6pt]
%
& \sum_{k \in \mathcal{K}} x_{j,k} \;\le\; 1
&& \forall\, j \in \mathcal{J} \setminus \mathcal{J}_C,
\tag{7} \label{eq:single} \\[6pt]
%
& y_j \;=\; 1
&& \forall\, j \in \mathcal{J}_M,
\tag{8} \label{eq:mandatory} \\[6pt]
%
& \sum_{a \in \mathcal{A}} f_{a,k}
  \;\le\;
  \mathcal{M} \sum_{a \in \delta^+(S)} f_{a,k}
&& \forall\, k \in \mathcal{K},
\tag{9} \label{eq:depot} \\[6pt]
%
& f_{a,k} \;\in\; \mathbb{Z}_{\ge 0}
&& \forall\, a \in \mathcal{A},\; k \in \mathcal{K},
\tag{10} \label{eq:dom-f} \\[4pt]
%
& x_{j,k} \;\in\; \{0,1\}
&& \forall\, j \in \mathcal{J} \setminus \mathcal{J}_C,\;
   k \in \mathcal{K},
\tag{11} \label{eq:dom-x} \\[4pt]
%
& y_j \;\in\; \{0,1\}
&& \forall\, j \in \mathcal{J} \setminus \mathcal{J}_C.
\tag{12} \label{eq:dom-y}
\end{align}

\section{Model Description}

The objective function~\eqref{eq:obj} seeks to maximise the total length of
cleaned streets minus a penalty for water waste.  The first term sums the
length of every street that is cleaned; the second term accounts for the
excess water used when a vehicle's capacity exceeds the street's
requirement, weighted by the penalty factor~$W$.  When $W = 0$ the problem
reduces to pure coverage maximisation.

Constraint~\eqref{eq:flow} enforces \emph{flow conservation} at every
junction for each vehicle: the total flow leaving a junction equals the
total flow entering it.  This ensures that each vehicle's route forms a
closed walk, so it can return to its starting point without discontinuities.

Constraint~\eqref{eq:time} limits each vehicle's total traversal time to
the shared time budget~$T$, guaranteeing that every vehicle returns to the
depot on time.

Constraint~\eqref{eq:traverse} requires that if vehicle~$k$ cleans
street~$j$, it must have traversed at least one arc of that street.  A
vehicle cannot clean a street it has not physically visited.

Constraints~\eqref{eq:link-fwd} and~\eqref{eq:link-rev} together link the
per-vehicle cleaning variables $x_{j,k}$ to the aggregate cleaned indicator
$y_j$.  Constraint~\eqref{eq:link-fwd} ensures that cleaning by any
vehicle activates the indicator ($x_{j,k} = 1 \Rightarrow y_j = 1$).
Constraint~\eqref{eq:link-rev} prevents the indicator from being active
unless at least one vehicle actually cleans the street ($y_j = 1
\Rightarrow \sum_k x_{j,k} \ge 1$).  Together they encode $y_j = 1
\Leftrightarrow \exists\, k : x_{j,k} = 1$.

Constraint~\eqref{eq:single} limits each street to be cleaned by at most
one vehicle.  Although the problem permits redundant cleaning, enforcing
single-vehicle cleaning eliminates unnecessary water waste, since a
street's length is scored only once regardless of how many vehicles clean
it.

Constraint~\eqref{eq:mandatory} ensures that every mandatory street is
cleaned.  A solution that leaves any mandatory street uncleaned is
infeasible.  If a mandatory street is provably unreachable within the time
budget (detected via a Dijkstra-based pre-check), its hard constraint is
relaxed to a soft one so the solver is not forced into infeasibility.

Constraint~\eqref{eq:depot} is a Big-M formulation that forces each
vehicle to include the depot in its route.  If a vehicle uses any arc
(left-hand side ${} > 0$), it must have at least one outgoing traversal from
the depot.  Combined with flow conservation~\eqref{eq:flow}, this ensures
the closed walk passes through~$S$.  When the depot has zero outgoing flow,
the right-hand side collapses to~$0$, prohibiting all arc usage --- i.e.\
an unused vehicle simply stays at the depot.

Constraints~\eqref{eq:dom-f}--\eqref{eq:dom-y} state the variable domains:
flow variables are non-negative integers, and cleaning and cleaned-indicator
variables are binary.

This model is a Mixed-Integer Linear Programming problem representing a
multi-vehicle arc routing problem with heterogeneous fleet, capacity
matching, time windows, and a bi-criteria objective balancing street
coverage against water efficiency, in the context of climate-aware urban
street cleaning.

\subsection{Subtour Elimination}

Flow conservation alone does not prevent \emph{subtours} --- disconnected
cycles that never visit the depot.  These are eliminated via an iterative
cutting-plane procedure.  After each solve, the weakly-connected components
of each vehicle's arc-flow multigraph are computed.  For every
component~$\mathcal{C}$ that does not contain the depot, the following
constraint is added:
%
\begin{equation}\tag{13}\label{eq:subtour}
\sum_{\substack{a \in \mathcal{A}:\\ \mathrm{tail}(a) \in \mathcal{C},\;
      \mathrm{head}(a) \in \mathcal{C}}} f_{a,k}
\;\le\;
\Biggl(\,1 + \!\!
  \sum_{\substack{a \in \mathcal{A}:\\ \mathrm{tail}(a) \in \mathcal{C},\;
        \mathrm{head}(a) \in \mathcal{C}}} U_a
\Biggr)
\cdot
\sum_{\substack{a \in \mathcal{A}:\\ \mathrm{tail}(a) \notin \mathcal{C},\;
      \mathrm{head}(a) \in \mathcal{C}}} f_{a,k}
\end{equation}
%
This constraint states that the total flow on arcs internal to~$\mathcal{C}$
can be positive only if there is positive flow on arcs entering~$\mathcal{C}$
from outside.  If no arcs enter the component, all internal flow is forced
to zero.  The model is then re-solved with the added cuts (warm-started
from the previous solution) and the process repeats until no subtours remain
or the time budget is exhausted.

\subsection{Post-Processing: Route Reconstruction}

After solving, each vehicle's arc-flow solution is converted into an
explicit junction sequence via Hierholzer's algorithm for Eulerian circuits.
If any connected component of a vehicle's multigraph is disconnected from
the depot, a shortest-path bridge (depot $\to$ component $\to$ depot) is
inserted when the remaining time budget permits; otherwise the component's
arcs are dropped.

\end{document}
